\emph{Setting.} As in Theorem \claimref{synthesis.rung-field-on-the-lift-locus}: $\mathcal{R}_{\C}$ has complex coordinates $(\alpha, \beta) \in \C^{2} \times \C^{2}$, $U(2)$ acts by $\alpha \mapsto X\alpha$, $\beta \mapsto \bar{X}\beta$, and by Proposition \claimref{standard-model.ten-transported-ambient-space-doublets-are} the $v$-coefficients carry a weak doublet of hypercharge $\tfrac12$ and the $w$-coefficients its conjugate. The Hermitian form is $\operatorname{Re}(c^{*}\cdot d)$ with $c^{*} = \theta(\bar{c})$, so $(\alpha, \beta)^{*} = (\bar{\beta}, \bar{\alpha})$ and $N(c) = \tfrac12(|\alpha|^{2} + |\beta|^{2})$. Let $D = \C^{2}$ with the same generators $T_{j} = \tfrac{\mathrm{i}}{2}\sigma_{j}$, $Y = \tfrac{\mathrm{i}}{2}\mathbf{1}$, viewed as a real representation of dimension four.

\emph{(1).} $D$ is irreducible as a real representation: $SU(2)$ acts transitively on the unit sphere of $D$, so a nonzero invariant real subspace contains a sphere's worth of directions and is all of $D$. Its real-linear equivariant endomorphisms commute with the central generator $Y$, which acts as $\tfrac{\mathrm{i}}{2}$, so they are complex-linear, and by Schur's lemma they are the complex scalars. The map $\beta \mapsto \bar{\beta}$ carries $(\C^{2}, \bar{X})$ to $(\C^{2}, X)$ equivariantly, so as a real representation $\mathcal{R}_{\C} \cong D \oplus D$, and the real-linear equivariant maps $D \to \mathcal{R}_{\C}$ are $\gamma \mapsto (p\gamma, q\bar{\gamma})$ with $p, q \in \C$, a real four-dimensional space; a nonzero one is injective because its kernel is an invariant subspace of the irreducible $D$. The script solves the intertwining equations for an unknown real $8 \times 4$ matrix against the four generators, by exact elimination, finds a solution space of real dimension four, and checks that it equals the span of the four maps $(p, q) = (1, 0), (\mathrm{i}, 0), (0, 1), (0, \mathrm{i})$, while the inequivariant map $\gamma \mapsto (\bar{\gamma}, 0)$ fails the equations. By complete reducibility for the compact group $U(2)$, an invariant real subspace of dimension four is isomorphic to $D$, so it is the image of such a map: $D_{0}$ for $q = 0$, $D_{\infty}$ for $p = 0$, and otherwise $\{(p\gamma, q\bar{\gamma})\} = \{\beta = a\bar{\alpha}\}$ with $a = q/\bar{p}$.

\emph{(2).} On $D_{0}$ and $D_{\infty}$ one of $\alpha$, $\beta$ vanishes, so $Q = \alpha^{\mathsf T}\beta = 0 \neq -H^{2}$. On $G_{a}$, $Q = \alpha^{\mathsf T}a\bar{\alpha} = a|\alpha|^{2}$, $N = \tfrac12(1 + |a|^{2})|\alpha|^{2}$ and $m = \tfrac12(1 - |a|^{2})|\alpha|^{2}$. So $Q = -H^{2}$ forces $a|\alpha|^{2} = -H^{2}$, that is $a = -r$ with $r > 0$ and $|\alpha|^{2} = H^{2}/r$; $U(2)$ preserves $G_{-r}$ and acts on $\alpha$ by $X$, and $SU(2)$ is transitive on the three-sphere $|\alpha|^{2} = H^{2}/r$, so $G_{-r} \cap L$ is one orbit. The lift point $c_{\lambda}$ has $\alpha = \tfrac{H^{2}}{\lambda}e_{4}$, $\beta = -\lambda e_{4}$, so it lies in $G_{a}$ exactly for $a = -\lambda^{2}/H^{2}$. By Theorem \claimref{synthesis.rung-field-on-the-lift-locus}(4) the frame of a charge-preserving point is fixed by $m$, here by $m/N = (1 - r^{2})/(1 + r^{2})$, a constant of the subspace. The script checks these identities exactly for five values of $a$ (real and complex) and three $\alpha$ each, and the membership of $c_{\lambda}$ for fifteen pairs $(H^{2}, \lambda)$.

\emph{(3).} $G_{-r}$ is a linear subspace, so its tangent space at $c_{\lambda}$ is $G_{-r}$ itself; the tangent space of $L$ is $\ker dQ$, $dQ(c)[x] = 2Q(c, x)$ (two real conditions). Since $G_{-r} \cap L$ is one orbit of dimension three (the stabiliser of $c_{\lambda}$ is a circle, Theorem \claimref{synthesis.rung-field-on-the-lift-locus}(1)), and the gauge tangent vectors $T_{j}c_{\lambda}$, $Yc_{\lambda}$ lie in both $G_{-r}$ and $\ker dQ$, the intersection is the gauge tangent provided it has dimension three, which follows directly: for $x = (\delta\alpha, -r\overline{\delta\alpha})$, $Q(c_{\lambda}, x) = \tfrac12(\alpha^{\mathsf T}\delta\beta + \beta^{\mathsf T}\delta\alpha) = -\tfrac{\lambda}{2}(\overline{\delta\alpha_{4}} + \delta\alpha_{4})$, using $rH^{2}/\lambda = \lambda$, which vanishes exactly when $\operatorname{Re}\delta\alpha_{4} = 0$, a real three-dimensional space. The rung unit $u_{\lambda} = (-\tfrac{H^{2}}{\lambda^{2}}e_{4}, -e_{4})$ is not in $G_{-r}$: there $\delta\beta = -r\overline{\delta\alpha} = +e_{4}$. The physical directions of $L$ at $c_{\lambda}$, the neutral line of $u_{\lambda}$ and the charged plane (Theorem \claimref{synthesis.rung-field-on-the-lift-locus}(3)), are by definition complementary to the gauge tangent inside $T_{c_{\lambda}}L$, so none lies in $G_{-r} \cap T_{c_{\lambda}}L$; and $c_{\lambda}$ lies in no other one-doublet subspace by (2). An invariant real subspace containing $c_{\lambda}$ contains its $U(2)$-span, which is $G_{-r}$ (four-dimensional, by (1) and (2)); adding $u_{\lambda} \notin G_{-r}$ gives a subspace of dimension greater than four, which by complete reducibility and $\mathcal{R}_{\C} \cong D \oplus D$ is $\mathcal{R}_{\C}$. The script computes, for fifteen pairs $(H^{2}, \lambda)$ by exact rank: $G_{-r} \cap \ker dQ$ of dimension three equal to the gauge tangent, $u_{\lambda}$ and the two charged directions (solved as the kernel of $dQ$, the gauge tangents and $u_{\lambda}$ for the Hermitian form) outside $G_{-r}$, the generated span of $\{c_{\lambda}, u_{\lambda}\}$ of real rank eight and that of $c_{\lambda}$ equal to $G_{-r}$.

\emph{(4).} Inside $G_{-r}$ the gauge tangent has dimension three, so its complement for $\operatorname{Re}(x^{*}\cdot y)$ is a line; $c_{\lambda}$ is orthogonal to the gauge tangent, since $\operatorname{Re}(c^{*}\cdot Tc) = 0$ for every anti-Hermitian generator $T$, and so spans it. Along it $dQ(c_{\lambda})[c_{\lambda}] = 2Q(c_{\lambda}) = -2H^{2}$, so it changes $Q$; by contrast $Q(c_{\lambda}, u_{\lambda}) = \tfrac12(\tfrac{H^{2}}{\lambda}\cdot(-1) + (-\lambda)(-\tfrac{H^{2}}{\lambda^{2}})) = 0$. The script finds the complement by exact elimination and checks both values. In the linear field the directions changing $Q$ are the neutral normal pair, of mass squared $2W' + 4N(k_{10} \pm r)$ (Theorem \claimref{synthesis.linear-rung-field-has-a-neutral-normal-pair-in-the-scalon-sum}; there $r$ denotes a coupling, not the graph parameter), which grows without bound as $k_{10} \to \infty$.

\emph{(5).} The radial line of $G_{-r}$ at $c_{\lambda}$ is that of $c_{\lambda}$, and $c_{\lambda}/\lambda - u_{\lambda} = (\tfrac{2H^{2}}{\lambda^{2}}e_{4}, 0)$ exactly (checked), while both have $\beta$-component $-e_{4}$; so the angle between the two directions is of order $H^{2}/\lambda^{2}$. Along $c = (1 + t)c_{\lambda}$, $Q(c) = (1 + t)^{2}Q(c_{\lambda})$, so $Q + H^{2} = -2H^{2}t + O(t^{2})$ and $k_{10}|Q + H^{2}|^{2} = 4k_{10}H^{4}t^{2} + O(t^{3})$; with the kinetic coefficient $N(c_{\lambda}) = \tfrac12(\lambda^{2} + H^{4}/\lambda^{2})$ for $t$ the added mass squared is of order $8k_{10}H^{4}/\lambda^{2}$ (computed, up to the normalisation convention of the kinetic term). At the stored frame $\lambda^{2}/H^{2} = e^{2\eta}$ is of order $10^{88}$, so both are negligible at any finite $k_{10}$.

\emph{Conclusion.} On the exact locus a one-doublet field is pure gauge, and its one physical direction is normal to the locus. The stored theory's neutral physical scalar is the rung unit (Theorem \claimref{synthesis.rung-field-d-term-action-and-tree-masses-on-the-lift-locus}(3)) and it carries a charged pair; by (3) a field whose restriction to $L$ carries both takes values in a subspace containing $c_{\lambda}$ and $u_{\lambda}$, which is all of $\mathcal{R}_{\C}$. This gives two doublets under three readings taken exactly: the null-lift constraint as a hard constraint, the Higgs exactly the rung unit, and a linear field. By (5) it does not follow at finite $k_{10}$, where a one-doublet field reproduces the rung-unit direction to relative order $H^{2}/\lambda^{2}$; and the hard limit $k_{10} \to \infty$ is a nonlinear limit, not a perturbative decoupling. The representation $\mathcal{R}_{\C}$ is that of the ten; it admits the one-doublet real forms $G_{a}$, among them $G_{1}$, fixed by the tower's conjugation ($\beta = \bar{\alpha}$, where $Q = |\alpha|^{2} > 0$), and $G_{-1}$, which meets $L$ only at the balanced frame (as in Theorem \claimref{synthesis.rung-unit-is-frame-boost}(5)); so the representation alone does not fix the count. Standard content: in two-doublet language $u_{\lambda}$ changes the ratio $|\alpha|/|\beta|$ of the two vacuum values at fixed $\alpha^{\mathsf T}\beta$, a direction a single doublet lacks. The result is transport-convenient: the summand, the locus and the frame are transport's, the argument is linear algebra.
